% Section 8.3, from lectures/L08/README.md: what you should be able to do afterwards, the
% questions to test yourself, and the next lecture.

\section{Review}
\label{c8:sec:review}

\needspace{6\baselineskip}
\subsection*{What you should be able to do now}
After this chapter, you should be able to:
\begin{itemize}
  \item Trace an interrupt from the device's line to your handler, naming the three numbers
    involved.
  \item Read \file{/proc/interrupts} and say what a stuck count and a racing count each indicate.
  \item Register a handler, and say what the \code{dev_id} argument is for on a shared line.
  \item Say when \code{IRQ_NONE} is the correct return value, and what breaks if you never return
    it.
  \item Acknowledge a level-triggered interrupt in the right order, and say what the wrong order
    loses.
  \item List what a hard interrupt handler may not do, and derive the list from one fact rather than
    memorising it.
  \item Choose between a workqueue and a threaded handler, and defend the choice.
\end{itemize}

\needspace{6\baselineskip}
\subsection*{Questions to test yourself}
\begin{enumerate}
  \item The device tree says interrupt 112 and \file{/proc/interrupts} says 47. Are they the same
    interrupt?
  \item Your handler runs once and the machine then stops responding. What did you forget?
  \item Why can a hard interrupt handler not call \code{copy_to_user}, given that it is only copying
    memory?
  \item On a shared line, your handler always returns \code{IRQ_HANDLED}. Describe the symptom on
    the \emph{other} device sharing that line.
  \item What is the difference between a tasklet and a workqueue, in terms of what each may do?
  \item You measure 9,847 interrupts where you predicted 10,000. Name two mechanisms that could
    account for the difference, and say how you would tell which it was.
\end{enumerate}

\needspace{6\baselineskip}
\subsection*{Where to look}
\begin{itemize}
  \item \file{tools/qa-dev/qa-dev.h} is the register map, and documents \code{IRQ_STATUS} as
    write-one-to-clear.
  \item The device's own model, \file{tools/qa-dev/qa-dev.c}, is readable and shows exactly when the
    line is raised and lowered. Reading the hardware's source is a luxury you will not have again.
\end{itemize}

\needspace{6\baselineskip}
\subsection*{Looking ahead}
Chapter~\ref{c9:ch} is about:
\begin{itemize}
  \item The other half: a reader that wants data which does not exist yet.
  \item Wait queues, and the lost wakeup that is a bug rather than bad luck.
  \item \code{poll}, and how one thread waits on your device and a socket at once.
  \item Jiffies, \code{ktime}, and the difference between delaying and sleeping.
\end{itemize}
